\documentclass[12pt,a4paper]{article}

\usepackage[utf8]{inputenc}
\usepackage[T1]{fontenc}
\usepackage[brazil]{babel}
\usepackage[a4paper,margin=2.cm]{geometry}
\usepackage{enumitem}
\setlist[enumerate]{itemsep=0.45em,topsep=0.35em}

\begin{document}

\begin{center}
  {\Large\bfseries PROVA SIMULADA -- NPOR -- VERSÃO 4}\par
  \vspace{0.4cm}
  Nome: \rule{9cm}{0.4pt}\par\vspace{0.3cm}
  Data: \rule{4cm}{0.4pt}\hfill Nota: \rule{3cm}{0.4pt}
\end{center}

\noindent\textbf{Duração sugerida: 90 minutos.} A prova oficial pode ter outro tempo e outros conteúdos; confirme as instruções da sua Organização Militar. Este simulado tem 13 questões objetivas e uma redação.

\medskip
\noindent\textbf{Instruções:} marque uma alternativa por questão. Na redação, escreva um texto dissertativo-argumentativo, com letra legível, respeitando o tema e o espaço disponível.

\subsection*{Português}
Leia o texto para responder às questões 1 e 2.

\begin{quote}
Áreas verdes nas cidades oferecem espaços para descanso e convivência. Além disso, árvores podem ajudar a reduzir a temperatura em ruas muito expostas ao sol. Para que esses benefícios permaneçam, o planejamento urbano precisa considerar a preservação e a manutenção desses locais.
\end{quote}

\begin{enumerate}
  \item Qual ideia é defendida pelo texto?
  \begin{enumerate}[label=\alph*)]
    \item As áreas verdes são úteis apenas para embelezar as cidades.
    \item A manutenção e a preservação são importantes para conservar os benefícios das áreas verdes.
    \item As árvores aumentam a temperatura das ruas.
    \item O planejamento urbano deve substituir as áreas verdes por construções.
  \end{enumerate}

  \item No contexto, a expressão ``Além disso'' indica:
  \begin{enumerate}[label=\alph*)]
    \item oposição.
    \item adição de uma informação.
    \item explicação de uma causa.
    \item conclusão de um argumento.
  \end{enumerate}

  \item Assinale a frase com concordância nominal adequada.
  \begin{enumerate}[label=\alph*)]
    \item As propostas foram analisadas pela equipe.
    \item As proposta foram analisadas pela equipe.
    \item Os moradores ficaram bastante satisfeitas.
    \item Segue anexas os documentos solicitados.
  \end{enumerate}

  \item Assinale a alternativa em que o uso da crase está correto.
  \begin{enumerate}[label=\alph*)]
    \item O representante dirigiu-se à comissão.
    \item A equipe começou à preparar o relatório.
    \item Os participantes voltaram à pé.
    \item A orientação foi dada à ele.
  \end{enumerate}
\end{enumerate}

\subsection*{Matemática e raciocínio lógico}
\begin{enumerate}[resume]
  \item Um capital de R\$ 1.000,00 foi aplicado a juros simples de 2\% ao mês durante 3 meses. Qual foi o valor dos juros?
  \begin{enumerate}[label=\alph*)]
    \item R\$ 20,00
    \item R\$ 60,00
    \item R\$ 120,00
    \item R\$ 600,00
  \end{enumerate}

  \item No sistema $2x+y=11$ e $x+y=7$, qual é o valor de $x$?
  \begin{enumerate}[label=\alph*)]
    \item 2
    \item 3
    \item 4
    \item 7
  \end{enumerate}

  \item Um triângulo tem base de 10 cm e altura de 6 cm. Qual é sua área?
  \begin{enumerate}[label=\alph*)]
    \item 16 cm$^2$
    \item 30 cm$^2$
    \item 60 cm$^2$
    \item 120 cm$^2$
  \end{enumerate}

  \item Uma caixa retangular mede 4 cm por 3 cm por 2 cm. Qual é seu volume?
  \begin{enumerate}[label=\alph*)]
    \item 9 cm$^3$
    \item 12 cm$^3$
    \item 24 cm$^3$
    \item 48 cm$^3$
  \end{enumerate}

  \item Complete a sequência: 1, 4, 9, 16, \ldots
  \begin{enumerate}[label=\alph*)]
    \item 20
    \item 24
    \item 25
    \item 32
  \end{enumerate}
\end{enumerate}

\subsection*{Conhecimentos gerais e institucionais}
\begin{enumerate}[resume]
  \item Em que ano foi assinada a Lei Áurea, que aboliu legalmente a escravidão no Brasil?
  \begin{enumerate}[label=\alph*)]
    \item 1822
    \item 1850
    \item 1888
    \item 1889
  \end{enumerate}

  \item Qual linha imaginária atravessa a região Norte do Brasil?
  \begin{enumerate}[label=\alph*)]
    \item Trópico de Capricórnio.
    \item Linha do Equador.
    \item Círculo Polar Ártico.
    \item Meridiano de Greenwich.
  \end{enumerate}

  \item Quais estados compõem a Região Sudeste?
  \begin{enumerate}[label=\alph*)]
    \item São Paulo, Rio de Janeiro, Minas Gerais e Espírito Santo.
    \item São Paulo, Paraná, Minas Gerais e Goiás.
    \item Rio de Janeiro, Bahia, Espírito Santo e Paraná.
    \item Minas Gerais, Mato Grosso, São Paulo e Rio de Janeiro.
  \end{enumerate}

  \item Qual posto vem imediatamente antes de Tenente-Coronel na hierarquia do oficialato?
  \begin{enumerate}[label=\alph*)]
    \item Capitão.
    \item Major.
    \item Primeiro-Tenente.
    \item Coronel.
  \end{enumerate}
\end{enumerate}

\newpage
\section*{Redação}
\noindent\textbf{Tema:} A importância da perseverança diante dos desafios pessoais e coletivos.

\medskip
\noindent Escreva um texto dissertativo-argumentativo, apresentando seu ponto de vista, argumentos que o sustentem e uma conclusão. Sugestão: de 15 a 25 linhas. Não é necessário escrever um título.

\vspace{0.4cm}
\noindent\rule{\textwidth}{0.3pt}\par\vspace{0.45cm}
\noindent\rule{\textwidth}{0.3pt}\par\vspace{0.5cm}
\noindent\rule{\textwidth}{0.3pt}\par\vspace{0.5cm}
\noindent\rule{\textwidth}{0.3pt}\par\vspace{0.5cm}
\noindent\rule{\textwidth}{0.3pt}\par\vspace{0.5cm}
\noindent\rule{\textwidth}{0.3pt}\par\vspace{0.5cm}
\noindent\rule{\textwidth}{0.3pt}\par\vspace{0.5cm}
\noindent\rule{\textwidth}{0.3pt}\par\vspace{0.5cm}
\noindent\rule{\textwidth}{0.3pt}\par\vspace{0.5cm}
\noindent\rule{\textwidth}{0.3pt}\par\vspace{0.5cm}
\noindent\rule{\textwidth}{0.3pt}\par\vspace{0.5cm}
\noindent\rule{\textwidth}{0.3pt}\par\vspace{0.5cm}
\noindent\rule{\textwidth}{0.3pt}\par\vspace{0.5cm}
\noindent\rule{\textwidth}{0.3pt}\par\vspace{0.5cm}
\noindent\rule{\textwidth}{0.3pt}\par\vspace{0.5cm}
\noindent\rule{\textwidth}{0.3pt}\par\vspace{0.5cm}
\noindent\rule{\textwidth}{0.3pt}\par\vspace{0.5cm}
\noindent\rule{\textwidth}{0.3pt}\par\vspace{0.5cm}
\noindent\rule{\textwidth}{0.3pt}\par\vspace{0.5cm}
\noindent\rule{\textwidth}{0.3pt}\par\vspace{0.5cm}
\noindent\rule{\textwidth}{0.3pt}\par\vspace{0.5cm}
\noindent\rule{\textwidth}{0.3pt}\par

\end{document}